\chapter{A finite-time physical access theorem}
\label{sec:physical-pulse}

An intervention library is scientifically useful only when its operations have
a physical realization.  The following result gives one such realization:
native probability current becomes an endpoint population response under a
selective kinetic intervention.  The result concerns finite-state classical
Markov dynamics.  It does not require reconstruction of the surrounding
network, infinitesimal perturbations, or an infinitesimal observation time.

Throughout this section, distributions are row vectors.  A finite row
generator $Q$ has nonnegative off-diagonal entries and $Q\mathbf 1=0$.
Let $\pi Q=0$, with $\pi$ a probability distribution, and let $i\ne j$.
Write $P'(t)=e^{tQ'}$ for the semigroup after a fixed intervention.  The
intervention remains in place until readout at the predetermined time $t$.

\begin{theorem}[Finite-amplitude endpoint response]
\label{thm:physical-endpoint-response}
Suppose $Q'$ is another valid finite generator and differs from $Q$ only in
the directed rates $i\to j$, $j\to i$, and their compensating diagonal
entries.  Define
\[
 \gamma=\pi_i(Q'_{ij}-Q_{ij})-\pi_j(Q'_{ji}-Q_{ji}).
\]
For every finite $t>0$,
\begin{align}
 [\pi P'(t)]_j-\pi_j&=\gamma H_{ij}'(t),\label{eq:physical-response}\\
 H_{ij}'(t)&=\int_0^t\bigl(P'_{jj}(s)-P'_{ij}(s)\bigr)\,ds\nonumber\\
 &=\int_0^t P'_{jj}(s)\,
       \Pr_i^{Q'}(\tau_j>t-s)\,ds>0,\label{eq:physical-hitting}
\end{align}
where $\tau_j$ is the first hitting time of $j$, starting at $i$.
Consequently the endpoint displacement has the sign of $\gamma$ for
every positive time, and vanishes at one such time if and only if
$\gamma=0$.  Irreducibility of either generator is unnecessary.
\end{theorem}

\begin{proof}
The assumptions give $\pi Q'=\gamma(e_j-e_i)$, where $e_k$ denotes the
row point mass at $k$.  Integrating
$\frac{d}{ds}\pi P'(s)=\pi Q'P'(s)$ proves
\eqref{eq:physical-response} and the first expression for $H$.
Let $F_j(t)=\int_0^tP'_{jj}(s)\,ds$.  Before $\tau_j$, a path starting
at $i$ spends no time in $j$.  The strong Markov property gives
\[
 \int_0^tP'_{ij}(s)\,ds
 =\mathbb E_i\!\left[
       \mathbf 1_{\{\tau_j\le t\}}F_j(t-\tau_j)\right].
\]
Interchanging nonnegative integrals yields
\eqref{eq:physical-hitting}.  Set $\lambda_k=-Q'_{kk}$.  A path
remaining in its initial state gives
$P'_{jj}(s)\ge e^{-\lambda_j s}$ and
$\Pr_i(\tau_j>u)\ge e^{-\lambda_i u}$.
Both factors are strictly positive at finite times, including when
$\tau_j=\infty$ has positive probability.  Their integral is positive.
\end{proof}

The integrand $P'_{jj}(s)-P'_{ij}(s)$ in the first line need not be
positive.  For example, a directed three-cycle has negative instantaneous
contrast at some times.  The sign theorem follows from the integrated
first-hitting representation, not from pointwise transition ordering.
It forbids crossing the \emph{native baseline}; it does not assert that
the displacement magnitude is monotone or cannot overshoot its eventual
stationary value.

\begin{corollary}[Local quantitative bounds]
\label{cor:physical-local-bounds}
If the perturbed exit rates of $i,j$ have upper bounds
$\Lambda_i,\Lambda_j$, then
\[
 0<L(t):=\int_0^t
 e^{-\Lambda_j s-\Lambda_i(t-s)}\,ds
 \le H'_{ij}(t)\le t.
\]
For a common upper bound $\Lambda$, this gives
$t e^{-\Lambda t}\le H'_{ij}(t)\le t$.
For unequal bounds,
$L(t)=(e^{-\Lambda_i t}-e^{-\Lambda_j t})/(\Lambda_j-\Lambda_i)$.
No rate bound on other states appears.
\end{corollary}

\section{Selective channels and native current}

A physical channel $e$ between $i$ and $j$ has forward and reverse rates
$a_e,b_e$.  Several physical channels may share the same state pair.
The native channel current is
$J_e=\pi_i a_e-\pi_j b_e$.  Suppose a calibrated intervention multiplies
\emph{both} rates of this channel by $r\in(0,1)$, leaving all other
channels unchanged.  Then $\gamma=-(1-r)J_e$, so that
\begin{equation}
 D_e(t):=\mathbb E_\pi[
 \mathbf 1_{\{X_0=j\}}-\mathbf 1_{\{X_t=j\}}]
 =(1-r)J_e H'_{ij}(t).
 \label{eq:physical-current-response}
\end{equation}
Thus one resolved endpoint population, compared before and after a
fixed kinetic step, identifies the orientation of the native channel
current in the ideal exact-law limit.  The stationary distribution
\emph{after} intervention generally changes and must not be substituted
for $\pi$ when defining that native current.

This access genuinely separates some passively indistinguishable
mechanisms.  On two observed states, the two-channel rate assignments
$(a_1,b_1,a_2,b_2)=(3,1,1,3)$, $(2,2,2,2)$, and $(1,3,3,1)$ all give
the same aggregate generator with rates $4$ in each direction.  Their
first-channel native currents are $1,0,-1$.  A factor-$1/2$ selective
intervention produces negative, zero, and positive endpoint displacements
according to \eqref{eq:physical-current-response}.  A common additive
state-level control applied to their identical aggregate generators
cannot make that distinction.

In an Arrhenius model a barrier increase can scale the two rates by
the same factor while preserving their ratio.  Whether a particular
chemical, transport, or electronic device provides this selective
control is an empirical calibration question.  A binding reporter,
an arbitrary inhibitor, or a change of a shared reservoir does not
automatically implement the required operation.

\section{Which coarse detectors preserve the theorem?}

\begin{corollary}[A sufficient entrance-gateway condition]
\label{cor:physical-gateway}
Let $B$ contain $j$ but not $i$.  If every directed path from $i$ into
$B$ under $Q'$ first visits $j$, then
\[
 H'_{i,j;B}(t):=\int_0^t
     \bigl(P'_j(X_s\in B)-P'_i(X_s\in B)\bigr)\,ds
 =\int_0^tP'_j(X_s\in B)\Pr_i(\tau_j>t-s)\,ds>0.
\]
Theorem~\ref{thm:physical-endpoint-response} therefore also holds for
the $B$-occupancy displacement, with this kernel.  The same local
lower bound applies since $P'_j(X_s\in B)\ge P'_{jj}(s)$.
\end{corollary}

The proof is the same first-hitting decomposition: occupation in $B$
before $\tau_j$ is zero.  The gateway property is checkable by removing
$j$ and testing whether $i$ can still reach any state in $B$.

A more general sufficient condition is a regenerative entrance law.
Let $\nu_A$ be supported outside $B$, let $\nu_B$ be supported inside
$B$, and suppose
$\pi Q'=\gamma(\nu_B-\nu_A)$.
If, starting from $\nu_A$, the entrance state at $\tau_B$ has law
$\nu_B$ independent of the entrance time, then the response of
$B$-occupancy is $\gamma$ times
\[
 \int_0^t\Pr_{\nu_B}^{Q'}(X_s\in B)
             \Pr_{\nu_A}^{Q'}(\tau_B>t-s)\,ds>0.
\]
Independence here means the joint entrance law factors, including
conditioning on any set of finite entrance times.  It does not follow
from generic coarse graining.  The strong Markov property after
entrance and this stated factorization prove the claim.

Without such an interface condition the sign can reverse classically.
Consider the four bidirectional channels
\[
 (i,j,a,b)=(0,1,2,1),\ (0,2,10,1),\ (1,3,10,1),\ (2,3,1,3).
\]
The native law is $(7,4,80,30)/121$.  Suppressing channel $0\leftrightarrow1$
by $1/2$ gives stationary law $(13,6,140,50)/209$.
For the coarse detector $B=\{1,2\}$, the initial displacement derivative
is $-5/121$ but the stationary displacement is $10/2299>0$.
Continuity therefore forces a finite-time crossing.  The resolved
singleton $\{1\}$ retains the sign theorem.  A measured crossing alone
is consequently not a device-independent witness of quantum dynamics.

\section{Calibration and finite independent records}

For classical noninvasive paired readout, define
$Z=\mathbf 1_{\{X_0=j\}}-\mathbf 1_{\{X_t=j\}}\in[-1,1]$.
From $N$ independently prepared records, Hoeffding's inequality gives
\[
 \Pr(|\overline Z-\mathbb EZ|\ge s)\le
           2\exp(-Ns^2/2).
\]
This classical paired observation is an additional access assumption;
an initial quantum measurement can disturb the prepared state.

Suppose the inferred response lies in $[d_-,d_+]$, a calibrated factor
lies in $[r_-,r_+]\subset(0,1)$, and both endpoint exit rates are at most
$\Lambda$.  If preparation is within total variation $\eta$ of native
stationarity, its contribution to paired-response bias is at most
$2\eta$.  Let $b$ bound this and all other propagated response biases.
Set
\[
 I_D=[d_--b,d_++b]\cap[-1,1],\qquad
 I_G=[(1-r_+)t e^{-\Lambda t},\ (1-r_-)t].
\]
An enclosing native-current interval is
\begin{equation}
 I_J=\left[
 \min_{d\in\partial I_D,\ g\in\partial I_G}\frac d g,
 \max_{d\in\partial I_D,\ g\in\partial I_G}\frac d g
 \right].\label{eq:physical-current-interval}
\end{equation}
The denominator is strictly positive.  The implementation replaces
the exponential by a rational Taylor-series enclosure.  If $I_J$
contains zero, the orientation is unresolved.  A joint calibration
failure allowance is added to the sampling allowance once; calibration
and sampling need not be independent for that union bound.

The fixed-pulse kernel interval generally leaves a nonzero magnitude
uncertainty even with infinitely many records.  One can calibrate the
kernel separately or shorten the pulse while increasing sample size.
The next bound makes the latter option explicit.

\begin{proposition}[A constructive magnitude-precision bound]
\label{prop:physical-magnitude}
Assume exact native-stationary preparation, an exactly calibrated
$a=1-r>0$, endpoint exit bound $\Lambda>0$, and
$|J_e|\le J_{\max}>0$.  If $\pi_j\le p_{\max}\le1$, then
\[
 \left|\frac{\mathbb EZ}{at}-J_e\right|
 \le J_{\max}\Lambda t,\qquad
 \operatorname{Var}(Z)\le
 V(t):=\min\{1,(2p_{\max}\Lambda+aJ_{\max})t\}.
\]
For target absolute error $\varepsilon>0$, choose
$t=\varepsilon/(2J_{\max}\Lambda)$ and
$s=at\varepsilon/2$.  The estimator $\overline Z/(at)$ has error at
most $\varepsilon$ with probability at least $1-\alpha$ whenever
\[
 N\ge\frac{2\log(2/\alpha)}{s^2}
       \left(V(t)+\frac{2s}{3}\right).
\]
For fixed nonzero model bounds this is a sufficient
$N=O(\varepsilon^{-3})$ record count.  It is not a minimax lower bound.
\end{proposition}

\begin{proof}
The response identity and $e^{-\Lambda t}\le H(t)/t\le1$ give the
bias bound, using $1-e^{-x}\le x$.  Since
$Z^2=\mathbf 1_{\{X_0=j\}}+\mathbf 1_{\{X_t=j\}}
       -2\mathbf 1_{\{X_0=j,X_t=j\}}$,
\[
 \mathbb EZ^2=2\pi_j(1-P'_{jj}(t))-D_e(t)
 \le2p_{\max}\Lambda t+aJ_{\max}t.
\]
Also $Z^2\le1$ and $|Z-\mathbb EZ|\le2$.
Bernstein's inequality bounds the sampling error by $s$ at the stated
count.  Dividing by $at$ contributes $\varepsilon/2$, and the
deterministic bias contributes the remaining $\varepsilon/2$.
\end{proof}

The precision bound is conditional.  A fixed unmodeled response bias
is amplified by $1/(at)$ as pulses shorten.  The appropriate finite
error method must retain this bias or abstain.  No rate of convergence
is claimed for arbitrary imperfect hardware.
If no tighter current cap is available, the same local calibration
supplies $J_{\max}=\Lambda/r$: the perturbed exit rates bound
$ra_e$ and $rb_e$, and
$|\pi_i a_e-\pi_j b_e|\le\max(a_e,b_e)$.  This bound can be
conservative but requires no reconstruction of the surrounding network.

\section{Physical resources and experimental acceptance}

Neither a kinetic suppression factor nor an equilibrium-preserving
rate ratio is a complete energy account.  An applied experiment must
include native fuel consumption, stationary preparation, resets,
selective control, detector operation, recording, and calibration.
If independent calibration supplies a per-record upper energy bound
$E_0+P_{\max}t$, Proposition~\ref{prop:physical-magnitude} implies
the corresponding \emph{sufficient} bound
$N(E_0+P_{\max}t)$.  For small $\varepsilon$, its preparation term
scales as $O(\varepsilon^{-3})$ and its active-pulse term as
$O(\varepsilon^{-2})$.  These conditional engineering bounds are not
universal thermodynamic precision laws.

The physical acceptance conditions are native stationary preparation;
a stable, time-homogeneous generator during each record; selective
two-endpoint control with independently bounded cross-effects; a
resolved endpoint or certified gateway detector; predetermined pulse
duration; and readout while the perturbed generator is still applied.
State-dependent stopping, feedback, releasing the perturbation before
readout, uncalibrated detector latency, and arbitrary hidden-state
aggregation require their own analysis.  Violating any of these
conditions can explain an observed failure of the sign law.

The implementation also supplies a stronger-access route based on
local transition counts.  Stopping at a fixed accumulated source-state
exposure makes a specified transition count Poisson with parameter
equal to rate times exposure.  It does not make counts conditionally
Poisson at fixed elapsed time.  Source exposure can take arbitrarily
long wall time; a deadline-limited acquisition must report incomplete
exposure rather than silently condition its statistical guarantee on
completion.
